\SecondArticle{13}
Der Beweis dieses Satzes wäre allerdings sehr leicht, wenn man annehmen dürfte, $Y$ lasse sich in lineare Faktoren zerlegen:
\[
(x-A)(x-B)(x-C)(x-D)\, .\, .\, .\, .\, .
\]
Dann wäre nämlich auch gewiss, dass $Z$ das Produkt aus allen $u-(A+B)x+AB$ und die Determinante der Funktion $Z$ das Produkt aus den Differenzen zwischen je zweien der Größen
\[
\begin{gathered}
(A+B)x-AB\\
(A+C)x-AC\\
(A+D)x-AD\\
\text{usw.}\\
(B+C)x-BC\\
(B+D)x-BD\\
\text{usw.}\\
(C+D)x-CD\\
\text{usw. usw.}
\end{gathered}
\]
ist. Dieses Produkt kann aber nicht identisch verschwinden, wenn nicht einer seiner Faktoren selbst identisch $=0$ wird. Daraus würde folgen, dass zwei der Größen $A,B,C$ usw. gleich sind und somit die Determinante $P$ des Polynoms $Y$ $=0$ wird, gegen die Voraussetzung.

\SourceLeaf{055}
Doch lassen wir diese Argumentation beiseite, die offenbar wie in Artikel 6 von einer Vorwegnahme des zu Beweisenden ausgeht, und gehen sogleich zu einem sicheren Beweis des Satzes von Artikel 12 über.

\SecondArticle{14}
Die Determinante der Funktion $\zeta$ ist das Produkt aus allen Differenzen zwischen je zweien der Ausdrücke $(a+b)x-ab$. Die Anzahl dieser Differenzen ist
\[
\frac12m(m-1)(\tfrac12m(m-1)-1)=\frac14(m+1)m(m-1)(m-2)
\]
Diese Zahl gibt also den Grad der Determinante der Funktion $\zeta$ hinsichtlich der Unbestimmten $x$ an. Die Determinante der Funktion $z$ hat denselben Grad. Dagegen kann die Determinante der Funktion $Z$ durchaus einen niedrigeren Grad haben, wenn einige Koeffizienten, von der höchsten Potenz von $x$ an, verschwinden. Wir haben nun zu beweisen, dass in der Determinante der Funktion $Z$ gewiss nicht \textit{alle} Koeffizienten verschwinden können.

Betrachten wir die Differenzen, deren Produkt die Determinante der Funktion $\zeta$ ist, näher, so finden wir: Ein Teil von ihnen, nämlich die Differenzen zwischen je zweien der Ausdrücke $(a+b)x-ab$ mit einem gemeinsamen Element, liefert
\[
\text{das Produkt aus allen}\quad(a-b)(x-c)
\]
Aus den übrigen, nämlich den Differenzen zwischen je zweien der Ausdrücke $(a+b)x-ab$ mit verschiedenen Elementen, entsteht dagegen
\[
\begin{gathered}\text{das Produkt aus allen}\quad(a+b-c-d)x-ab+cd,\\\text{mit ausgeschlossenen Wiederholungen.}\end{gathered}
\]
Das erste Produkt enthält offenbar jeden Faktor $a-b$ genau $m-2$ Mal und jeden Faktor $x-c$ genau $(m-1)(m-2)$ Mal. Daraus schließen wir leicht, dass dieses Produkt
\[
=\pi^{m-2}\upsilon^{(m-1)(m-2)}
\]
ist. Bezeichnen wir also das zweite Produkt mit $\rho$, so ist die Determinante der Funktion $\zeta$
\[
=\pi^{m-2}\upsilon^{(m-1)(m-2)}\rho
\]
Bezeichnen wir ferner mit $r$ diejenige Funktion der Unbestimmten $x,l',l'',l'''$ usw., die durch die Substitutionen $l'=\lambda',l''=\lambda'',l'''=\lambda'''$ usw. in $\rho$ übergeht, und mit $R$
\SourceLeaf{056}
die Funktion allein von $x$, in die $r$ durch die Substitutionen $l'=L',l''=L''$, $l'''=L'''$ usw. übergeht, so ist offenbar die Determinante der Funktion $z$
\[
=p^{m-2}y^{(m-1)(m-2)}r
\]
die Determinante der Funktion $Z$ dagegen
\[
=P^{m-2}Y^{(m-1)(m-2)}R
\]
Da nach Voraussetzung $P$ nicht $=0$ ist, kommt es nun darauf an, zu beweisen, dass $R$ gewiss nicht identisch verschwinden kann.

\SecondArticle{15}
Dazu führen wir noch eine weitere Unbestimmte $w$ ein und betrachten das Produkt aus allen
\[
(a+b-c-d)w+(a-c)(a-d)
\]
mit ausgeschlossenen Wiederholungen. Da es $a,b,c$ usw. symmetrisch enthält, lässt es sich als Polynom in den Unbestimmten $w,\lambda',\lambda'',\lambda'''$ usw. darstellen. Wir bezeichnen diese Funktion mit $f(w,\lambda',\lambda'',\lambda'''\text{ usw.})$. Die Anzahl jener Faktoren $(a+b-c-d)w+(a-c)(a-d)$ ist
\[
=\frac12m(m-1)(m-2)(m-3)
\]
woraus wir leicht schließen, dass
\[
f(0,\lambda',\lambda'',\lambda'''\text{ usw.})=\pi^{(m-2)(m-3)}
\]
und daher auch
\[
f(0,l',l'',l'''\text{ usw.})=p^{(m-2)(m-3)}
\]
sowie
\[
f(0,L',L'',L'''\text{ usw.})=P^{(m-2)(m-3)}
\]
gilt. Das Polynom $f(w,\allowbreak L',\allowbreak L'',\allowbreak L'''\text{ usw.})$ hat im Allgemeinen den Grad
\[
\frac12m(m-1)(m-2)(m-3)
\]
in besonderen Fällen kann es aber durchaus einen niedrigeren Grad haben, wenn etwa einige Koeffizienten, von der höchsten Potenz von $w$ an, verschwinden. Unmöglich ist dagegen, dass die ganze Funktion identisch $=0$ ist, denn
\SourceLeaf{057}
die eben gefundene Gleichung zeigt, dass wenigstens ihr letztes Glied nicht verschwindet. Wir nehmen an, das höchste Glied des Polynoms $f(w,\allowbreak L',\allowbreak L'',\allowbreak L'''\text{ usw.})$ mit nicht verschwindendem Koeffizienten sei $Nw^\nu$. Setzen wir nun $w=x-a$, so ist $f(x-a,L',L'',L'''\text{ usw.})$ offenbar ein Polynom in den Unbestimmten $x,a$ oder, was dasselbe ist, ein Polynom in $x$ mit Koeffizienten, die von der Unbestimmten $a$ abhängen. Sein höchstes Glied ist dabei jedoch $Nx^\nu$ und hat somit einen bestimmten, von $a$ unabhängigen Koeffizienten, der nicht $=0$ ist. Ebenso sind $f(x-b,L',L'',L'''\text{ usw.}),f(x-c,L',L'',L'''\text{ usw.})$ Polynome in der Unbestimmten $x$, deren höchstes Glied jeweils $Nx^\nu$ ist, während die Koeffizienten der folgenden Glieder jeweils von $b,c$ usw. abhängen. Das Produkt aus den $m$ Faktoren
\[
\begin{gathered}
f(x-a,L',L'',L'''\text{ usw.})\\
f(x-b,L',L'',L'''\text{ usw.})\\
f(x-c,L',L'',L'''\text{ usw.})\\
\text{usw.}
\end{gathered}
\]
ist daher ein Polynom in $x$, dessen höchstes Glied $N^mx^{m\nu}$ ist, während die Koeffizienten der folgenden Glieder von den Unbestimmten $a,b,c$ usw. abhängen.

Wir betrachten nun weiter das Produkt aus den folgenden $m$ Faktoren:
\[
\begin{gathered}
f(x-a,l',l'',l'''\text{ usw.})\\
f(x-b,l',l'',l'''\text{ usw.})\\
f(x-c,l',l'',l'''\text{ usw.})\\
\text{usw.}
\end{gathered}
\]
Da es eine Funktion der Unbestimmten $x,a,b,c$ usw., $l',l'',l'''$ usw. ist, und zwar hinsichtlich $a,b,c$ usw. symmetrisch, lässt es sich als Funktion der Unbestimmten $x,\lambda',\lambda'',\lambda'''$ usw., $l',l'',l'''$ usw. darstellen, die wir mit
\[
\varphi(x,\lambda',\lambda'',\lambda'''\text{ usw.},\ l',l'',l'''\text{ usw.})
\]
bezeichnen. Somit ist
\[
\varphi(x,\lambda',\lambda'',\lambda'''\text{ usw.},\ \lambda',\lambda'',\lambda'''\text{ usw.})
\]
das Produkt aus den Faktoren
\[
\begin{gathered}
f(x-a,\lambda',\lambda'',\lambda'''\text{ usw.})\\
f(x-b,\lambda',\lambda'',\lambda'''\text{ usw.})\\
f(x-c,\lambda',\lambda'',\lambda'''\text{ usw.})\\
\text{usw.}
\end{gathered}
\]
\SourceLeaf{058}
und daher identisch durch $\rho$ teilbar, da man leicht sieht, dass jeder Faktor von $\rho$ in einem jener Faktoren enthalten ist. Wir setzen also
\[
\begin{aligned}
&\varphi(x,\lambda',\lambda'',\lambda'''\text{ usw.},\ \lambda',\lambda'',\lambda'''\text{ usw.})\\
&=\rho\psi(x,\lambda',\lambda'',\lambda'''\text{ usw.})
\end{aligned}
\]
wobei das Zeichen $\psi$ ein Polynom bezeichnet. Daraus folgt aber leicht, dass auch identisch
\[
\begin{aligned}
&\varphi(x,\allowbreak L',\allowbreak L'',\allowbreak L'''\text{ usw.},\allowbreak \ L',\allowbreak L'',\allowbreak L'''\text{ usw.})\\
&=R\psi(x,L',L'',L'''\text{ usw.})
\end{aligned}
\]
gilt. Oben haben wir jedoch bewiesen, dass das Produkt aus den Faktoren
\[
\begin{gathered}
f(x-a,L',L'',L'''\text{ usw.})\\
f(x-b,L',L'',L'''\text{ usw.})\\
f(x-c,L',L'',L'''\text{ usw.})\\
\text{usw.}
\end{gathered}
\]
das $=\varphi(x,\lambda',\lambda'',\lambda'''\text{ usw.},\ L',L'',L'''\text{ usw.})$ ist, das höchste Glied $N^mx^{m\nu}$ hat. Dasselbe höchste Glied hat daher auch die Funktion $\varphi(x,\allowbreak L',\allowbreak L'',\allowbreak L'''\text{ usw.},\allowbreak $
$L',L'',L'''\text{ usw.})$, und sie ist somit gewiss nicht identisch $=0$. Folglich kann auch $R$ nicht identisch $=0$ sein und daher ebenso wenig die Determinante der Funktion $Z$. Q. E. D.

